%\vspace{-0.3cm}
\section{Related Work} \label{sec:related_works}
%\vspace{-0.3cm}
System identification is a broadly studied topic~\citep {ASTROMSystemID, schoukens, SCHON201139, ziemann2022single, ziemann2022learning}. However, system identification from the perspective of experiment design for nonlinear systems is much less understood~\citep{chiuso2019system}. Most methods formulate the identification task through the maximization of intrinsic rewards.
 %Intrinsic rewards depend on the RL agent's model of the system and encourage exploration by suggesting policies that visit transitions where the learned model has less ``information''. 
Common choices of intrinsic rewards are (\emph{i}) model prediction error or ``Curiosity''~\citep{schmidhuber1991possibility, pathak2017curiosity}, (\emph{ii}) novelty of transitions~\citep{stadie2015incentivizing}, and  (\emph{iii}) diversity of skills~\citep{eysenbach2018diversity}.
%. Curiosity approaches encourage collecting data where the learned model makes the most prediction error. Thus, the model error is used as a proxy for lack of information. However, these methods ...
%Alternatively, methods with state visitation count as the intrinsic reward encourage visiting states with low visitation frequency, that is the states that have been less frequently observed by the agent. For finite MDPs 

A popular choice for intrinsic rewards is mutual information or entropy~\citep{learning_and_control_using_gps, active_learning_gp, max, pathak2019self, sekar2020planning}. \cite{learning_and_control_using_gps} propose an
approach to maximize the information gain greedily wrt the immediate next transition, i.e., one-step greedy, whereas \cite{active_learning_gp} consider planning full trajectories.
\cite{max, pathak2019self, sekar2020planning} and \cite{Sancaktaretal22} consider general Bayesian models, such as BNNs, to represent a probabilistic distribution for the learned model. 
\cite{max} propose using the information gain of the model with respect to observed transition as the intrinsic reward. %Accordingly, the algorithm suggests policies that visit transitions which give more information for the learned model. 
To this end, they learn an ensemble of Gaussian neural networks and represent the distribution over models with a Gaussian mixture model (GMM). A similar approach is also proposed in \cite{pathak2019self, sekar2020planning, Sancaktaretal22}. The main difference between \cite{max} and \cite{pathak2019self} lies in how they represent mutual information. Moreover, \cite{pathak2019self} use the model's epistemic uncertainty, that is the disagreement between the ensemble models as an intrinsic reward. \cite{sekar2020planning} link the model disagreement (epistemic uncertainty) reward to maximizing mutual information and demonstrate state-of-the-art performance on several high-dimensional tasks. 
Similarly, \cite{Sancaktaretal22}, use the disagreement in predicted trajectories of an ensemble of neural networks to direct exploration. Since trajectories can diverge due to many factors beyond just the model epistemic uncertainty, e.g., aleatoric noise, this approach is restricted to deterministic systems and susceptible to systems with unstable equilibria. 
Our approach is the most similar to \cite{pathak2019self, sekar2020planning} since we also propose the model epistemic uncertainty as the intrinsic reward for planning. However, we thoroughly and theoretically motivate this choice of reward from a Bayesian experiment design perspective. Furthermore,
we induce additional exploration in \opacex through our optimistic planner and rigorously study the theoretical properties of the proposed methods. On the contrary, most of the prior work discussed above either uses the mean planner (\textsc{Mean-AE}) or does not discuss the planner thoroughly or provide any theoretical results. In general, theoretical guarantees for active exploration algorithms are rather immature~\citep{chakraborty2023steering, wagenmaker2023optimal} and mostly restrictive to a small class of systems~\citep{simchowitz2018learning, tarbouriech2020active, wagenmaker2020active, mania2020active}. %In this work, we derive an upper bound on the information gain for \opacex for well-calibrated models. In addition, for the GP case, we show convergence for a rich class of RKHS, i.e.,  we show that information gain goes to zero. 
To the best of our knowledge, we are the first to give convergence guarantees for a rich class of nonlinear systems.

\looseness=-1
While our work focuses on the active learning of dynamics, there are numerous works that study exploration in the context of reward-free RL \citep{jin2020reward, kaufmann2021adaptive, wagenmaker2022reward, chen2022statistical}. However, most methods in this setting give guarantees for special classes of MDPs~\citep{jin2020reward, kaufmann2021adaptive, wagenmaker2022reward, qiu2021reward, chen2022statistical} and result in practical algorithms. On the contrary, we focus on solely learning the dynamics. While a good dynamics model may be used for zero-shot planning, it also exhibits more relevant knowledge about the system such as its stability or sensitivity to external effects. Furthermore, our proposed method is not only theoretically sound but also practical. 
%Another approach is proposed by
%\cite{Sancaktaretal22}, where the disagreement in predicted trajectories of an ensemble of neural networks is used to direct exploration. Moreover, instead of using intrinsic rewards for policy optimization, \cite{Sancaktaretal22} use the iCEM black box optimizer~\citep{iCem}
%to suggest actions that maximize disagreement in model trajectories. However, the approach is restricted to deterministic systems and susceptible to systems with unstable equilibria. In particular, for systems with unstable equilibria, small uncertainties around the equilibrium, can lead to drastic disagreement in trajectories. Accordingly, the algorithm might suggest policies that visit the unstable equilibrium even if the dynamics are accurate around that region. 